\section{Recopie de texte} 
%\label{sec:5:1:recopie} 

Dans le programme {\idt CroixGrecque\_1} ci-dessous on retrouve quatre fois la même
séquence de commandes pour dessiner les quatre
bran\-ches identiques de la croix et tourner. Chaque branche se décompose en deux cotés égaux
de lon\-gueur 3 et une extrémité de lon\-gueur 1 (l'unité
étant la longueur d'un côté d'un carreau  du quadrillage).

\vspace{\baselineskip}
\vspace*{3cm}%
\includevspacePostScript[-\baselineskip]{CroixGrecque}{0.95}{0.95}
\vspace*{-3cm}%
\label{prog:CroixGrecque_1}
\input CroixGrecque_1.tex

Cette situation où l'on trouve la même séquence de commandes du \robot\ 
plusieurs fois de suite dans un programme est très fréquente car les 
dessins géométriques sont souvent composés de morceaux identiques. Il est alors
très fastidieux de taper plusieurs fois le même texte. On peut l'éviter
en utilisant les possibilités offertes par l'ordinateur de copier un texte
déjà enregistré dans sa mémoire et de le recopier ailleurs dans le 
programme. La manière de procéder est expliquée pour les compatibles 
{\sc IBM PC} à la figure \ref{fig:PC-copier-coller} 
(page \pageref{fig:PC-copier-coller})
et pour le Macintosh à la figure \ref{fig:mac-copier-coller} 
(page \pageref{fig:mac-copier-coller}).
\input COPIER-COLLER-PC.tex
\input COPIER-COLLER-MAC.tex

Faire beaucoup de recopies à la main est long. De plus on imagine bien que 
dans un programme de plusieurs pages, le lecteur a du mal à retrouver les 
parties de texte identiques. Ce phénomème est illustré dans le
programme {\idt Nombre888\_1} ci-dessous. La difficulté serait
évidemment encore plus grande si les parties de texte identiques se
trouvaient sur des pages différentes.

\vspace{\baselineskip}
\vspace*{2cm}%
\includevspacePostScript[-\baselineskip]{Nombre_888}{0.95}{0.95}
\vspace*{-2cm}%
\input Nombre888_1.tex

\section{Exemples de procédures} 
%\label{sec:5:2:ex-procedures} 

Plutôt que de faire des recopies de textes à la main soi-même, on peut 
demander au compilateur \pascal\ de le faire à  notre place en utilisant 
des {\em procédures}\index{Procédure}. Commençons par réécrire le programme 
{\idt CroixGrecque\_1} (page \pageref{prog:CroixGrecque_1}) en uti\-li\-sant des procédures pour mettre en 
évidence la décomposition de la croix grecque en 
quatre bran\-ches identiques, chacune d'elles se décomposant en deux 
cotés égaux de lon\-gueur 3 et une extrémité de lon\-gueur 1.

\newpage
\vspace*{3cm}%
\includePostScript{CroixGrecque}{0.95}{0.95}
\vspace*{-3cm}%
\input CroixGrecque_2.tex

\noindent L'emploi d'une procédure comprend~:
\begin{list}{$\scriptscriptstyle\bullet$}{
}
\item la {\em déclaration de procédure}\index{Déclaration de procédure} qui définit une séquence 
d'instructions (appelée {\em corps de la procédure})\index{Corps de procédure} et lui donne un 
nom (appelé {\em nom de la procédure})\index{Identificateur}$\:$;
\item l'{\em appel de procédure}\index{Appel de procédure} qui indique que la séquence 
d'instructions cons\-ti\-tuant le corps de la procédure doit être recopiée 
à la place du nom de la procédure.
\end{list}

\noindent Pour le lecteur du programme, l'utilisation de procédures est 
préférable aux recopies de textes. A chaque fois qu'il voit le
nom d'une procédure, il comprend qu'il s'agit de la copie d'un texte
qu'il peut consulter dans la déclaration de procédure. S'il rencontre
plusieurs appels d'une même procédure, il n'a pas à faire d'effort 
pour comprendre immédiatement qu'il s'agit du même texte qui est
recopié. Le programme est donc plus clair. N'est-ce pas le cas du
programme {\idt Nombre888\_2} ci-dessous qui est plus facilement 
compréhensible que le programme {\idt Nombre888\_1}$\:$? 

\vspace{\baselineskip}
\vspace*{2cm}%
\includevspacePostScript[-\baselineskip]{Nombre_888}{0.95}{0.95}
\vspace*{-2cm}%
\input Nombre888_2.tex

\section{Déclaration de procédures} 
%\label{sec:5:3:decl-procedures} 

Les déclarations des procédures\index{Déclaration de procédure} d'un programme se placent entre {\mcl uses} et
{\mcl begin}. La structure d'un programme \pascal\ est donc la suivante~:

\begin{center}
\begin{tabular}{@{}|p{12.5cm}|@{}} \hline
{\em Programme \pascal\ :}\\ \hline \hline
\begin{minipage}[t]{12.5cm}
{\medskip
\mbox{\blc\ \ \ {\mcl program}\hspace{0.3em}\blanc\hspace{0.3em}\idt NomDuProgramme$;$}\par 
\mbox{\blc\ \ \ \ \mcl uses\idt\ \ldots\ $;$}\par 
\mbox{\blc\ \ \ \ \small\rm  \ldots\ Déclarations des procédures du programme \ldots}\par 
\mbox{\blc\ \ \ \mcl begin}\par 
\mbox{\blc\ \ \ \ \small\rm  \ldots\ Instructions du programme%
\footnote{\ Y compris les appels des procédures du programme.} (séparées par des `$;$') \ldots}\par 
\mbox{\blc\ \ \ \mcl end\idt$.$}\par 
\medskip}
\end{minipage}\\ \hline 
\end{tabular}
\end{center}


\noindent En \pascal, une déclaration de procédure a la forme suivante
(nous utiliserons des formes plus élaborées dans la suite)~:

\begin{center}
\begin{tabular}{@{}|p{12.5cm}|@{}} \hline
{\em Déclaration de procédure~:}\\ \hline \hline
\begin{minipage}[t]{12.5cm}
{\medskip
\mbox{\blc\ \ \ \mcl procedure\blc\hspace{0.3em}\blanc\hspace{0.3em}\idt NomDeLaProcedure$;$}\par 
\mbox{\blc\ \ \ \mcl begin}\par 
\mbox{\blc\ \ \ \ \small\rm \ldots\ Instructions du corps de la procédure%
\footnote{\ Y compris les appels de procédures précédemment déclarées.}
 (séparées par des `$;$') \ldots}\par 
\mbox{\blc\ \ \ \mcl end\idt$;$}\par 
\medskip}
\end{minipage}\\ \hline 
\end{tabular}
\end{center}

Le {\idt  NomDeLaProcedure} est choisi par le programmeur
à sa convenance, de façon à évoquer ce
à quoi sert cette procédure.
On notera qu'il faut un point-virgule `{\idt $;$}' après le {\idt  NomDeLaProcedure},
qu'il n'en faut pas après le {\mcl begin}, mais qu'il en faut un après le {\mcl end}
qui termine le corps de la procédure.


\section{Identificateurs} 
%\label{sec:7:2:Ident} 
Le nom d'un programme ou d'une procédure s'appelle un 
{\em identificateur}\index{Identificateur} en \pascal. En voici 
une définition précise à laquelle il
faudra se référer à chaque fois que nous parlerons d'iden\-ti\-fi\-ca\-teurs~:

\begin{center}
\begin{tabular}{@{}|p{12.5cm}|@{}} \hline
{\em Identificateur~:}\\ \hline \hline
\begin{minipage}[t]{12.5cm}
{\medskip
Un {\em identificateur} est une séquence de lettres 
minuscules ou majuscules non accentuées, de chiffres ou
de soulignés `\_' qui commence obligatoirement  par une lettre et qui est différente des
mots anglais réservés par \pascal\ comme {\mcl program}, {\mcl uses}, {\mcl procedure}, 
{\mcl begin}, {\mcl end}, {\mcl for}, {\mcl if}, etc.
\medskip}
\end{minipage}\\ \hline 
\end{tabular}
\end{center}

Par exemple {\idt Cote}, {\idt CroixGrecque} ou {\idt Croix\_Grecque} sont des
identificateurs. Pour rendre lisible un identificateur formé de plusieurs mots, on
met une majuscule au début de chaque mot, comme dans {\idt CroixGrecque} ou
{\idt EpaisseurCrayon}.

\section{Appel de procédures} 
%\label{sec:5:4:app-procedures} 

\index{Appel de procédure}Pour appeler une procédure, il suffit d'écrire son nom. Le {\idt NomDeLaProcedure}
est alors remplacé par le corps de la procédure\index{Corps de procédure} c'est-à-dire par la liste
d'instructions comprise entre le {\mcl begin} et le {\mcl end} dans sa déclaration.

On notera que si l'on déclare une procédure sans jamais l'appeler, elle 
ne sert absolument à rien. 

On peut appeler une procédure dans le programme ou dans une autre procédure.
Par exemple dans le programme {\idt CroixGrecque\_2}, on appelle 
la procédure {\idt Croix} qui appelle la procédure {\idt Branche} qui elle-même
appelle les procédures {\idt Cote} et {\idt Extremite}. 

Il faut impérativement respecter la règle qu'une procédure doit 
être
dé\-cla\-rée avant d'être ap\-pe\-lée. Supposons par exemple que l'on ait
déclaré une pro\-cé\-du\-re sous le nom {\idt NomProc}. Supposons également
que par suite d'une faute de frappe, on ait appelé cette procédure
sous le nom {\idt NomProg}. Le compilateur nous indique alors une erreur
`identificateur inconnu'. Ne sachant pas que nous avons écrit par erreur 
{\idt NomProg} au lieu de {\idt NomProc}, le compilateur suppose que la 
procédure {\idt NomProc} a été déclarée mais n'a jamais été utilisée.
C'est inutile mais autorisé. Il constate également que la procédure
{\idt NomProg} a été appelée mais n'a pas été déclarée. C'est interdit.
Retenant la faute la plus grave, le compilateur indique de manière
lapidaire que le nom {\idt NomProg} aurait
dû figurer dans une déclaration de procédure.%
\footnote{\ Nous parlons souvent des programmes de manière 
anthropomorphique, c'est-à-dire en leur attribuant des caractères 
propres à l'homme. En fait nous pensons aux programmeurs qui les
ont écrits. Par exemple, l'ordinateur ne se trompe pas$\:$; c'est son 
constructeur qui s'est trompé ou le programmeur qui a fait une erreur 
logique dans son programme.} 
On obtient la même erreur `identificateur inconnu' si l'on appelle une procédure
dont la déclaration est placée plus loin dans le programme.

\section{Commandes du robot définies par des procédures} 
%\label{sec:5:5:robot-cmd-procs}

Certaines commandes du \robot\ sont équivalentes à une séquence
de commandes élémentaires et peuvent donc s'exprimer avec des procédures.
Par exemple, pour pivoter d'un quart de tour, il suffit de pivoter deux fois
de suite d'un huitième de tour. On a donc~:
\medskip

\pagebreak[3] 
\vspace{\neglskip}\mbox{\blc  \idt  {\mcl procedure}\ pqtd$;$}\par 
\vspace{\neglskip}\mbox{\blc  \ \idt  {\cmt  \{\ Fait\ \underline{p}ivoter\ le\ robot\ sur\ place\ d{'}un\ \underline{q}uart\ de\ \underline{t}our\ à\ \underline{d}roite\ \}}}\par 
\vspace{\neglskip}\mbox{\blc  \idt  {\mcl begin}}\par 
\vspace{\neglskip}\mbox{\blc  \ \idt  phtd$;$\ phtd$;$\ }\par 
\vspace{\neglskip}\mbox{\blc  \idt  {\mcl end}$;$}\par 
\medskip

\pagebreak[3] 
\vspace{\neglskip}\mbox{\blc  \idt  {\mcl procedure}\ pqtg$;$\ }\par 
\vspace{\neglskip}\mbox{\blc  \ \idt  {\cmt  \{\ Fait\ \underline{p}ivoter\ le\ robot\ sur\ place\ d{'}un\ \underline{q}uart\ de\ \underline{t}our\ à\ \underline{g}auche\ \}}}\par 
\vspace{\neglskip}\mbox{\blc  \idt  {\mcl begin}}\par 
\vspace{\neglskip}\mbox{\blc  \ \idt  phtg$;$\ phtg$;$\ }\par 
\mbox{\blc  \idt  {\mcl end}$;$}\par 
\medskip

Voici de nouvelles commandes du \robot\ qui peuvent s'exprimer à l'aide de
procédures en utilisant des commandes que nous connaissons déjà~:
\medskip

\pagebreak[3] 
\vspace{\neglskip}\mbox{\blc  \idt  {\mcl procedure}\ p3htd$;$}\par 
\vspace{\neglskip}\mbox{\blc  \ \idt  {\cmt  \{\ Fait\ \underline{p}ivoter\ le\ robot\ sur\ place\ de\ \underline{3}\ \underline{h}uitièmes\ de\ \underline{t}our\ à\ \underline{d}roite\ \}}}\par 
\vspace{\neglskip}\mbox{\blc  \idt  {\mcl begin}}\par 
\vspace{\neglskip}\mbox{\blc  \ \idt  pqtd$;$\ phtd$;$\ }\par 
\vspace{\neglskip}\mbox{\blc  \idt  {\mcl end}$;$}\par 
\medskip

\pagebreak[3] 
\vspace{\neglskip}\mbox{\blc  \idt  {\mcl procedure}\ p3htg$;$}\par 
\vspace{\neglskip}\mbox{\blc  \ \idt  {\cmt  \{\ Fait\ \underline{p}ivoter\ le\ robot\ sur\ place\ de\ \underline{3}\ \underline{h}uitièmes\ de\ \underline{t}our\ à\ \underline{g}auche\ \}}}\par 
\vspace{\neglskip}\mbox{\blc  \idt  {\mcl begin}}\par 
\vspace{\neglskip}\mbox{\blc  \ \idt  pqtg$;$\ phtg$;$\ }\par 
\vspace{\neglskip}\mbox{\blc  \idt  {\mcl end}$;$}\par 
\medskip

\pagebreak[3] 
\vspace{\neglskip}\mbox{\blc  \idt  {\mcl procedure}\ pdt$;$}\par 
\vspace{\neglskip}\mbox{\blc  \ \idt  {\cmt  \{\ Fait\ \underline{p}ivoter\ le\ robot\ sur\ place\ d{'}un\ \underline{d}emi\ \underline{t}our\ \}}}\par 
\vspace{\neglskip}\mbox{\blc  \idt  {\mcl begin}}\par 
\vspace{\neglskip}\mbox{\blc  \ \idt  pqtd$;$\ pqtd$;$\ }\par 
\vspace{\neglskip}\mbox{\blc  \idt  {\mcl end}$;$}\par 
\medskip

\noindent\rule{137mm}{0.1mm} %-------------------------------------------------------
\begin{exercices}{
Écrire des programmes \pascal\ comportant des pro\-cé\-du\-res pour 
réaliser les croix,
%de la figure \ref{fig:Croix}, 
l'étoile astroïdale, la loupe ou le double carré cranté
%de la figure \ref{fig:EtLoCa} 
ci-dessous.
\noindent $\Box$}\end{exercices}
\begin{figure}[htbp]
 \vspace{-\baselineskip}
 \includevspacePostScript[1.7cm]{Croix}{0.95}{0.95}
 \caption{Croix potencée, double, papale, tréflée et de Saint André}
 %\label{fig:Croix}
\end{figure}
\begin{figure}[htbp]
 \includevspacePostScript[5.1cm]{EtoileLoupeCarre}{0.95}{0.95}
 \caption{Étoile astroïdale, loupe et double carré cranté}
 %\label{fig:EtLoCa}
\end{figure}

\clearpage

\begin{corrige}{
Nous donnons un exemple de programme, les
autres se trouvant corrigés sur la disquette d'accompagnement.
\input EtoileAstroidale.tex
\noindent $\Box$}\end{corrige}

%%%%%%%%%%%%%%%%%%%%%%%%%%%%%%%%%%%%%%%%%%%%%%%%%%%%%%%%%%%%%%%%%%%%%%%%%%%
\noindent \framebox[\textwidth]{
\begin{minipage}{\myparboxwidth}
\setlength{\parindent}{2em}
\small \sf
\vspace{\abovedisplayskip}
\section*{Addenda}
La notion de procédure\index{Procédure} introduite par Von Neumann et Turing
fût utilisée en 1949
sous le nom de {\em routine} sur l'EDSAC ({\em Electronic Delay
Storage Automatic Computer}) généralement considéré 
comme le premier ordinateur à programme enregistré ayant 
fonctionné. Cette machine fût construite dès 1947 à l'université
de Cambridge en Angleterre par l'équipe de Wilkes\footnote{\ Maurice
Wilkes, l'un des inventeurs du radar.} et exécuta ses premiers programmes
en Mai 1949 \cite{levy1}.

Les procédures illustrent le concept d'{\em abstraction}\index{Abstraction}~: les 
propriétés d'un objet sont isolées afin de
les utiliser ou de les étudier à part. Plus précisément, une 
procédure donne un nom à un algorithme de construction d'une partie
d'un dessin. Dans le corps de la procédure, 
on s'intéresse à la façon de construire le dessin. Dans les appels 
on s'intéresse uniquement au dessin terminé, tout en désirant complètement ignorer comment 
le dessin est obtenu. On sépare donc l'objet et la façon de le 
construire, le quoi du comment. 

\vspace{\abovedisplayskip}
\end{minipage}
}
%%%%%%%%%%%%%%%%%%%%%%%%%%%%%%%%%%%%%%%%%%%%%%%%%%%%%%%%%%%%%%%%%%%%%%%%%%%%
%%%%%%%%%%%%%%%%%%%%%%%%%%%%%%%%%%%%%%%%%%%%%%%%%%%%%%%%%%%%%%%%%%%%%%%%%%%
\noindent \framebox[\textwidth]{
\begin{minipage}{\myparboxwidth}
\setlength{\parindent}{2em}
\small \sf
\vspace{\abovedisplayskip}
Nous trouvons utile de décrire le dessin réalisé par une procédure
dans un commentaire 
placé au début de la procédure que l'on appelle la
 {\em spécification}\index{Spécification} de la procédure. On remarquera que pour qu'une procédure 
puisse être utilisée sans jamais consulter son corps, il faut que la 
spécification de la procédure décrive quel dessin elle fait mais 
également quelles sont les positions et orientations
du \robot\ avant et après la construction du dessin. Il se dégage ainsi
une notion d'{\em invariant}\index{Invariant} sur laquelle nous reviendrons dans les chapitres
suivants sur l'itération. 

Les procédures sont à la base des méthodes 
de structuration des très grands programmes~: pour être 
compréhensibles, les programmes de plusieurs milliers de lignes 
doivent être découpés en procédures ne dépassant pas 
quelques pages. Les procédures constituent donc un moyen élégant pour 
segmenter logiquement un programme. Cependant le concept de procédure
est insuffisant pour structurer des programmes de plusieurs millions de
lignes. On utilise alors la notion de {\em module} (on dit également {\em unité})
et de {\em bibliothèque} pour regrouper des ensembles de procédures.

Des procédures regroupées dans une
{\em bibliothèque} (en Turbo Pascal on dirait {\em unité}) peuvent être 
fournies aux programmeurs pour les aider dans leurs tâches. Pour utiliser 
une procédure, le programmeur doit connaître son nom et sa 
spécification sans avoir besoin
de connaître son fonctionnement interne. La clause {\mcl uses} au début
du programme a pour but d'indiquer au compilateur quelles sont les 
bibliothèques utilisées dans ce programme. Les commmandes du \robot\ sont
des procédures regroupées dans des bibliothèques.

L'analyse de dessins géométriques pour en dégager une structure que
l'on peut décrire algorithmiquement à l'aide de procédures est très
formatrice pour l'esprit. En ce sens l'apprentissage de la programmation
peut jouer pour les enfants un rôle formateur analogue à celui que 
jouait autrefois la géométrie euclidienne.

\vspace{\abovedisplayskip}
\end{minipage}
}
%%%%%%%%%%%%%%%%%%%%%%%%%%%%%%%%%%%%%%%%%%%%%%%%%%%%%%%%%%%%%%%%%%%%%%%%%%%%
\newpage
\thispagestyle{empty}